\section{Верификация и численная точность}
\label{sec:verification}

\subsection{Постановка проверок MMS}
\draft{Используется исходное уравнение~\eqref{eq:heat} с переменными $x,z,t$.
В методе изготовленных решений~\cite{salari2000} задаётся гладкое поле $T_*$,
из которого вычисляются источник, начальные и граничные данные.
Предусмотреть единую серию на одном слое с неподвижными и движущимися границами;
положительная теплопроводность линейно меняется по координатам:}
\begin{gather}
 \xi=\frac{x}{W},\qquad
 \eta=\frac{z-z_{\mathrm{surf}}(t)}{h(t)},\qquad
 \lambda_* = \lambda_0(1+\alpha\xi+\beta\eta),\qquad
 \rho c=C_0,\\
 T_*(x,z,t)=T_{\mathrm{base}}+
 B\psi(t)\cos(\pi\xi)\sin\left(\frac{\pi\eta}{2}\right),\\
 A_* = C_0\left(\frac{\partial T_*}{\partial t}
             +\boldsymbol v\cdot\nabla T_*\right)
             -\nabla\cdot(\lambda_*\nabla T_*).
\end{gather}
\draft{Здесь $\lambda_0,C_0>0$, $\alpha,\beta>0$;
скорость материала определяется движением границ слоя.
Для стационарной проверки $\psi=1$, геометрия неподвижна.
Для нестационарной проверки $\psi=\exp(-t/t_{\mathrm{final}})$;
границы либо неподвижны, либо линейно меняют положение.
Начальное поле задаётся по $T_*(x,z,0)$. Это проверки численного оператора
с изготовленным источником, а не расчёт геологического набора.}
\draft{Для сеточных проверок используется ошибка:}
\begin{equation}
 E_T=\left(\frac{\sum_P|P|\bigl(T_P-T_*(x_P,z_P,t)\bigr)^2}
                    {\sum_P|P|}\right)^{1/2}.
\end{equation}
\draft{В нестационарном расчёте она измеряется при $t=t_{\mathrm{final}}$.
Описанная единая серия с переменной теплопроводностью и движущимся слоем
ещё не выполнена; ниже сохранены имеющиеся результаты до обновления рисунков.}

\subsection{Пространственная точность и перенос}
\draft{Рисунок~\ref{fig:spatial} показывает имеющуюся пространственную проверку
при постоянной теплопроводности и перенос между сетками.
Последние оценки порядка --- 2,14 для MPFA,
1,89 для NMPFA и 1,98 для переноса.}
\reportfigure{05_spatial_convergence.png}
  {Ошибка температуры на гладком изготовленном решении и при переносе поля
  между сетками. Пунктиром показано квадратичное уменьшение ошибки.}
  {fig:spatial}

\subsection{Контакты материалов и баланс потоков}
\draft{На рис.~\ref{fig:contacts} отдельно проверены горизонтальный
и наклонный контакты с контрастом теплопроводности 100:1 и 1000:1.
Ошибка горизонтального теста достигает уровня округления; наклонный
контакт требует сгущения. Максимальный относительный дефект баланса
в серии --- $6{,}8\cdot10^{-14}$.}
\reportfigure{06_material_contacts.png}
  {Проверка сопряжения материалов: ошибка T, нормированная на перепад
  $\Delta T=20\,{}^\circ$C, при сгущении сетки.}
  {fig:contacts}

\subsection{Совместное сгущение по пространству и времени}
\draft{После отдельной стационарной проверки пространственного порядка
проверить нестационарную задачу MMS на неподвижном и движущемся слое.
При втором порядке по пространству и первом по времени одновременно
уменьшать $\Delta x,\Delta z$ вдвое, а $\Delta t$ --- вчетверо.
В асимптотическом режиме ожидается уменьшение общей ошибки в четыре раза:}
\begin{equation}
 E_T=\mathcal O(\Delta x^2+\Delta z^2+\Delta t),\qquad
 E_T\!\left(\frac{\Delta x}{2},\frac{\Delta z}{2},\frac{\Delta t}{4}\right)
 \approx\frac14 E_T(\Delta x,\Delta z,\Delta t).
\end{equation}
\draft{Рисунок~\ref{fig:time} содержит прежние вспомогательные проверки
временного интегратора и будет заменён результатами этой серии.}
\reportfigure{07_temporal_convergence.png}
  {Предварительные проверки временной точности; рисунок подлежит замене.
  Контрольная линия соответствует первому порядку.}
  {fig:time}

\subsection{Минимальная высота ячеек}
\draft{Влияние сжатых ячеек проверено через тепловое сопротивление
прослойки мощностью $\delta h=N\varepsilon$:
$\Delta T=q\delta h(1/\lambda_{\mathrm{new}}-1/\lambda_{\mathrm{old}})$
ниже прослойки. При $N=40$, $\varepsilon=0{,}05$ м в стационарном тесте
получено не более $0{,}075\,{}^\circ$C и 0,140\% относительного отличия Ro.
Рисунок~\ref{fig:minimum} сопоставляет эту оценку с чувствительностью полного расчёта.}
\reportfigure{08_minimum_cell_height.png}
  {Влияние суммарной мощности сжатых ячеек на T и Ro в стационарном тесте;
  справа --- изменение результата полного расчёта при другом $\varepsilon$.}
  {fig:minimum}

\subsection{Точность EASY\%Ro}
\draft{На рис.~\ref{fig:easy-accuracy} заданные истории нагрева и охлаждения
сопоставлены с независимым квадратурным интегрированием, проверенным удвоением
числа интервалов. Максимальная относительная ошибка Ro при табулировании
на этих историях --- $1{,}35\cdot10^{-6}\%$.}
\reportfigure{09_easy_ro_accuracy.png}
  {Четыре температурные истории, рассчитанная зрелость и относительная
  ошибка Ro при использовании таблицы относительно независимой квадратуры.}
  {fig:easy-accuracy}

\subsection{Синтетическая обратная задача}
\draft{Синтетическая обратная задача на рис.~\ref{fig:synthetic}:
известный поток по форме data1, геометрия обратного набора, 550 наблюдений
без добавленного шума. Все наблюдения согласованы с допуском, но ранний
поток отличается до 22,68\%; после продолжения подбора --- до 20,30\%.
Точность наблюдений и восстановление всей истории потока --- разные критерии.}
\reportfigure{10_synthetic_inverse.png}
  {Заданный и восстановленные потоки; справа --- распределения относительных
  отклонений синтетических T и Ro до и после продолжения подбора.}
  {fig:synthetic}

